% Person C: Enrollment and Reporting (implementation, usage and results)
\section{Enrollment and Reporting}

\subsection{Overview}
The enrollment module lets a student request courses, lets the administrator approve or reject the
requests, and lets both sides see the results. It is split into the same layers as the rest of the
project: \texttt{EnrollmentDAO} holds all SQL, \texttt{EnrollmentService} holds the rules and the role
checks, and the Swing panels (\texttt{BrowsePanel}, \texttt{StudentCoursesPanel}, \texttt{TimetablePanel},
\texttt{AdminEnrollmentsPanel}) only display data and call the service.

\subsection{Enrollment life cycle}
A request starts as \texttt{PENDING}. The administrator moves it to \texttt{APPROVED} or
\texttt{REJECTED}. A student can drop a course, which sets \texttt{DROPPED} (the row is kept for history).
Requesting the course again re-opens the old row, which respects the rule that a student cannot enroll in
the same course twice (\texttt{UNIQUE (student\_id, course\_id)}).

\subsection{Rules enforced}
\begin{itemize}
  \item Only an active student can request a course, and only an active (not archived) course accepts requests.
  \item Duplicate requests are rejected.
  \item The administrator can open or close the \emph{enrollment period}. While it is closed, students
        cannot request courses; approvals and drops still work. The setting is stored in a small
        \texttt{settings} table created on first use, so \texttt{schema.sql} did not change.
  \item Capacity is checked when a request is approved, so the number of approved students never exceeds
        the course capacity.
  \item Students see and change only their own enrollments; every service method checks the role first.
\end{itemize}

\subsection{Capacity and CGPA-based allocation}
When more students request a course than it has seats, the administrator can run automatic allocation.
Pending students are ranked by CGPA (highest first, ties broken by roll number) and approved until the
course is full. The remaining requests stay pending. The CGPA is computed by \texttt{CgpaCalculator} from
published grades, using the course credits and the assignment weights. Until grades are published, every
CGPA is $0.0$, so allocation falls back to the tie-break order.

\subsection{Timetable clashes and prerequisites}
Our final schema is for a single semester and gives every course a \emph{slot}. A partial unique index
allows only one active course per slot, so two active courses can never clash in time. Because of this
rule, no separate clash check is needed in the enrollment code, and the schema has no prerequisite table.

\subsection{Reports and CSV export}
The service provides four reports for the administrator: enrollment counts per course, course roster,
grade averages per course, and student transcript (percentage and letter grade per course, with the CGPA
in the title). Only \emph{published} grades are used. \texttt{CsvExporter} writes any report table to a CSV
file. The same reports are available from the command line through \texttt{python/reports.py}, which can
also draw a grade-average chart and a bell curve for one course.

\subsection{Usage}
\begin{enumerate}
  \item Run \texttt{bash scripts/setup.sh} to create the database, then \texttt{bash scripts/run.sh}.
  \item As administrator, open \emph{Enrollments} and use the \emph{Enrollment: OPEN/CLOSED} button.
  \item As a student, open \emph{Browse Courses}, search, and press \emph{Enroll}. A clear message is shown
        for errors such as a full course or a closed enrollment period.
  \item As administrator, filter pending requests by course and approve or reject them, or run the allocation.
  \item Students see their courses under \emph{My Courses} and the weekly grid under \emph{Timetable}.
\end{enumerate}

\subsection{Results}
The module is covered by 22 JUnit tests (\texttt{EnrollmentServiceTest}, \texttt{CsvExporterTest} and
\texttt{EnrollmentPeriodAndReportsTest}); the whole project passes 109 tests. The profiling results for
the enrollment queries on a database of 10{,}006 students and 50{,}040 enrollments are given in the
profiling section.
